\documentclass[openany,fontset=fandol]{ctexbook}
\usepackage{amsmath,amssymb}
\usepackage[unicode,hidelinks]{hyperref}
\newcommand{\note}[1]{\footnote{译注：#1}}
\begin{document}
\title{代数几何基础第五卷预备笔记\\引理13.3：证明的归约}
\author{据用户上传英译编辑稿译出}
\date{2026年11月17日}
\maketitle
\chapter*{范围与校订状态}
本批连续译出底本第51页从“It remains only to prove Lemma 13.3”起的一整段，止于“an integer $\geq2$”。为便于阅读，按论证和编写批注分段；没有省去原括注。下一处仍在第51页：“This is a regret to No. 9”。

本批是引理证明的归约开头，不是引理13.3或定理13.2完整证明。原稿的字母错置据同段上下文整理，原式保存在脚注；一般截面的基域及维数记号疑点保留待核。旧批的8.8、9.5缺文等未解决。本批尚未并入原合订本。
\chapter*{§13 初等态射与M. Artin定理（续）}
现在只剩下证明引理13.3。按通常的做法，只需证明经过$x$、余维为$n-1$的一般线性子簇（记为$L$）具有性质a)、b)、c)。\note{此处“一般”对应原文generic；后文用相应参数空间的一般点说明其含义，不是任意线性子簇。经过$x$这一条件在本段明确出现，与上一批指出的引理陈述遗漏相呼应。}

性质a)，以及b)、c)中的维数部分，可直接由2.3（已经在第12节重新审视并作修正）得到：如第8节中已用过的一次论证那样，将它应用于参数化经过$x$的直线的射影空间，以及$Z$经以$x$为中心的锥投影在该射影空间中的像。\note{底本此处印“passing through $X$”，但同句投影中心及整段指定点均为$x$，故订为$x$。“[dans ledit]”是底本保留的法语，指前述空间；正文译出这一指代。原文称维数结论直接来自2.3，本批不另补写其未展开的推导。}

［此外，将用这种方法所得的某些结果明确写出，或许也有益处：这些结果涉及由满足经过一个固定线性子簇这一条件的线性子簇所截出的线性截面。可放在正文中，也可另设一节。］\note{方括号内是原稿的编写建议，不是本译稿新增的证明或已完成的补充章节。}

对于b)、c)中的光滑性，由于a)，可以分别用$X\setminus Z$和$Y\setminus Z$代替$X$和$Y$；它们都是光滑的。于是问题归结为证明以下断言。

\newpage
设$f:X\to P$是一个未分歧态射，$X$在$k$上光滑，且$x\in P(k)$不属于$X$中任何维数小于$m$的分量的像（不可约分量？）。\note{括号中的疑问“irreducible component?”来自底本，照录而不替作者作定论。底本后半句印“let $X\in P(k)$”；因$X$已是态射的定义域，此处应为点$x$，正文据前后文整理，且将该点的说明前移以免悬空。}
令$\eta$为$\operatorname{Grass}^{m}P$中由经过$x$的余维$m$线性子簇$L$组成的子Grassmann簇的一般点。则相应的一般线性截面在$k$上光滑：至少在$k$的特征为零时如此；在另一种情形中，则须将$f$换成$Q_n f$，其中$n$为不小于2的整数。\note{底本写“$X^{(m)}$ is smooth over $k$”，未在这里印一般点下标；后续第52页明确写$X^{(m)}_\eta$。正文用“一般线性截面”表达其上下文所指，而不把整个关联族都断言为光滑。底域措辞仍照原文保留：一般纤维自然定义在$\kappa(\eta)$上，“在$k$上光滑”的字面表述与有限表示条件的关系需要核对，不能当作已核实的概形态射断言。底本此处确实印$Q_n f$，紧接的下一段改用$\Phi_n f$；此项符号不一致保留待核，不静默将二者认作同一映射。这里的$n$又用作该替换的整数指标，与段首维数记号重名，照录并提示。}
\end{document}
